\section{Algoritmet kryesore dhe zbatimi në Python}
Ky kapitull përmban tri familje problemesh. Secila kërkon një kriter ndalimi të lidhur me mbetjen, jo vetëm me ndryshimin ndërmjet dy iteracioneve.

\subsection{Përgjysmimi dhe Njutoni i mbrojtur}
\begin{lstlisting}[language=Python,caption={Përgjysmimi me kontroll të intervalit.}]
def pergjysmimi(f, a, b, tol=1e-12, max_iter=200):
    fa, fb = f(a), f(b)
    if fa * fb > 0:
        raise ValueError("Intervali nuk kufizon nje ndryshim shenje.")
    for _ in range(max_iter):
        c = 0.5 * (a + b)
        fc = f(c)
        if abs(fc) < tol or 0.5*(b-a) < tol:
            return c
        if fa * fc <= 0:
            b, fb = c, fc
        else:
            a, fa = c, fc
    raise RuntimeError("Nuk u arrit konvergjenca.")
\end{lstlisting}

\begin{lstlisting}[language=Python,caption={Njutoni me kontroll të derivatit.}]
def njutoni(f, df, x, tol=1e-12, max_iter=50):
    for _ in range(max_iter):
        fx, dfx = f(x), df(x)
        if abs(fx) < tol:
            return x
        if abs(dfx) < np.sqrt(np.finfo(float).eps):
            raise ArithmeticError("Derivat shume i vogel.")
        hapi = -fx / dfx
        x = x + hapi
        if abs(hapi) < tol * (1 + abs(x)):
            return x
    raise RuntimeError("Nuk u arrit konvergjenca.")
\end{lstlisting}
Në një zbatim robust, hapi i Njutonit pranohet vetëm nëse qëndron brenda intervalit dhe zvogëlon mbetjen; përndryshe kryhet një përgjysmim.

\subsection{Zgjidhja e sistemit linear}
\begin{lstlisting}[language=Python,caption={Zgjidhja dhe diagnostika pa formuar inversin.}]
A = np.array([[4., -1., 0.], [-1., 4., -1.], [0., -1., 3.]])
b = np.array([15., 10., 10.])
x = np.linalg.solve(A, b)
r = b - A @ x
print("x =", x)
print("mbetja relative =", np.linalg.norm(r)/np.linalg.norm(b))
print("kappa_2 =", np.linalg.cond(A))
\end{lstlisting}
Operatori \texttt{solve} përdor faktorizime të përshtatshme. Shkrimi \texttt{inv(A) @ b} shmanget.

\subsection{Metoda fuqi}
\begin{lstlisting}[language=Python,caption={Vlera vetjake dominante dhe kuocienti Rayleigh.}]
def metoda_fuqi(A, tol=1e-11, max_iter=10_000):
    x = np.ones(A.shape[0])
    x /= np.linalg.norm(x)
    lam_old = 0.0
    for k in range(max_iter):
        y = A @ x
        x = y / np.linalg.norm(y)
        lam = x @ A @ x
        if abs(lam - lam_old) < tol * (1 + abs(lam)):
            return lam, x, k + 1
        lam_old = lam
    raise RuntimeError("Metoda fuqi nuk konvergoi.")
\end{lstlisting}
Kontrolli përfundimtar është norma $\|A\vect x-\lambda\vect x\|$. Shenja e vektorit vetjak nuk ka kuptim fizik më vete: $\vect x$ dhe $-\vect x$ përfaqësojnë të njëjtin mod.

\subsection{Lidhja me modet normale}
Për $K\vect u=\omega^2M\vect u$, përdoret \texttt{scipy.linalg.eigh(K, M)} kur $K$ dhe $M$ janë simetrike dhe $M$ pozitivisht e përcaktuar. Frekuencat renditen, normalizohen dhe kontrollohet ortogonaliteti $U^TMU=I$.

Rendi i paraqitjes së metodave të rrënjëve, sistemeve dhe vlerave vetjake ndjek traditën e teksteve shqip \cite{hanelli,hoxha,datja}; termat e algjebrës lineare numerike kontrollohen edhe kundrejt punimeve të Departamentit të Matematikës së Aplikuar të UT-së \cite{zeqiri2020}.
